\documentclass[10pt]{article}
\usepackage{geometry}
\geometry{a4paper,scale=0.8}
\usepackage{amsmath,amsthm,amssymb,amscd}
\usepackage{latexsym}
\usepackage{indentfirst}
\usepackage{subfigure}
\usepackage{graphicx}
\usepackage{extarrows}
\usepackage{bm}
\usepackage{multirow}
\usepackage{algorithm}
%\usepackage{algorithmic}
\usepackage{algpseudocode}
\usepackage{setspace}
\usepackage{stackengine}
\usepackage{braket}
\usepackage{mathtools}
\usepackage{physics}
%\usepackage{commath}
\parskip 1em

\newtheorem{definition}{Definition}
\newtheorem{theorem}{Theorem}
\newtheorem{lemma}{Lemma}
\newtheorem{remark}{Remark}
\newtheorem{condition}{Condition}
\newtheorem{operation}{Operation}
\newtheorem{proposition}{Proposition}
\newtheorem{corollary}{Corollary}
\newtheorem{assumption}{Assumption}
\newtheorem{problem}{Problem}
\renewcommand{\algorithmicrequire}{\textbf{input:}}
\renewcommand{\algorithmicensure}{\textbf{initialization:}}
\usepackage{color,xcolor}
\DeclareMathOperator*{\argmax}{arg\,max}
\DeclareMathOperator*{\argsup}{arg\,sup}
\usepackage{authblk}
\usepackage{fncylab}
\labelformat{enumii}{\theenumi(#1)}
%\date{\today}



%\bibliographystyle{unsrt}
%\bibliographystyle{plain}
\linespread{2}


\begin{document}

\title{
  \bf 
}

\author[1]{Jingwei Li}
\affil[1]{Shanghai Center for Systems Biomedicine, Shanghai Jiao Tong University, Shanghai 200240, China, (Email:ljw2017@sjtu.edu.cn).}


%\pacs{}
%
%\keywords{}

\maketitle

\begin{definition}\label{DefData}
Define $M>0$ be an integer. Define $\mathcal{L}\coloneqq\set{\qty(d_m)_{m=0}^{M'-1} \mid 0<M'\le M, d_{M'-1}\neq 0, \forall 0\le m < M'-1, d_m > 0}$. Define $\mathcal{T}$ such that $\forall N>0$, $\set{d^{(n)} \mid 0\le n < N}\in \mathcal{T}$ iff $\forall 0\le n_1,n_2,n_3 < N$, $d^{(n_1)}\le d^{(n_2)}\le d^{(n_3)}$ and $d^{(n_1)},d^{(n_3)}>0$ implies that $d^{(n_2)}>0$.
\end{definition}

\begin{remark}
$\mathcal{T}$ in Definition \ref{DefData} can be defined equivalently as follows. $\forall \mathcal{D}\in\mathcal{T}$, $\mathcal{D}$ does not contain zero and the absolute values of negative elements in $\mathcal{D}$ do not split (even not strictly) those of positive elements. In other words, the positive values in $\mathcal{D}$ are consecutive when sorting the elements in absolute values.
\end{remark}

\begin{algorithm}
\caption{IMT}
\label{IMT}
\begin{algorithmic}[1]
\Require $D\coloneqq\set{d_m^{(n)}\neq 0\le n < N,0 \mid 0\le m < M_n}$.
\State $D_0\gets \set{d_0^{(n)} \mid 0\le n < N, M_n > 0}$.
\State Get the partition $\mathcal{E}_0$ of $D_0$ based on the equivalence relationship $\sim_0$ over $D_0$ such that $\forall d_0^n,d_0^{n'}\in D_0$, $d_0^n\sim_0 d_0^{n'}$ iff $d_0^n d_0^{n'}>0$ and $\forall d_0^{n''}\in D_0$, $\min\set{\abs{d_0^n},\abs{d_0^{n'}}}<\abs{d_0^{n''}}<\max\set{\abs{d_0^n},\abs{d_0^{n'}}}$ implies that $d_0^n d_0^{n''}>0$.
\State $F_0\gets \mathcal{P}(\Phi)\setminus\set{\bot}$.
\Repeat
\State $I\gets I+1$.
\State $F_I\gets \qty(\qty(\bigcup_{\psi\in F_{I-1}}D_{\mathbb{M}(\psi)}^A(\psi))\setminus\mathcal{U}\qty(\bigcup_{i=0}^{I-1}F_i))\land F_0$.
\Until{$F_I=\emptyset$}
\State \Return $\mathcal{F}(\Phi)\gets\set{F_i \mid 0\le i< I}$.
\end{algorithmic}
\end{algorithm}


%\centerline{\bf \bfseries  ---------------------------------------------------------}

\bibliographystyle{unsrt}
\bibliography{reference}

\end{document}
